\documentclass[aspectratio=169,11pt]{beamer}

\usepackage{fontspec}
\setsansfont{DejaVu Sans}
\setmonofont{DejaVu Sans Mono}[Scale=0.86]
\usepackage{booktabs}
\usepackage{array}
\usepackage{listings}
\usepackage{xcolor}
\usepackage{xurl}
\usepackage{tikz}
\usetikzlibrary{arrows.meta,positioning,shapes.geometric}

\usetheme{Madrid}
\usecolortheme{seahorse}
\setbeamertemplate{navigation symbols}{}
\setbeamertemplate{footline}[frame number]

\definecolor{sqlblue}{RGB}{28,74,135}
\definecolor{sqlgreen}{RGB}{24,117,74}
\definecolor{sqlgray}{RGB}{245,247,250}
\definecolor{sqlred}{RGB}{166,45,45}

\lstdefinestyle{sqlstyle}{
  language=SQL,
  basicstyle=\ttfamily\scriptsize,
  keywordstyle=\color{sqlblue}\bfseries,
  commentstyle=\color{sqlgreen},
  stringstyle=\color{sqlred},
  showstringspaces=false,
  breaklines=true,
  columns=fullflexible,
  backgroundcolor=\color{sqlgray},
  frame=single,
  rulecolor=\color{gray!35},
  xleftmargin=1mm,
  xrightmargin=1mm,
  aboveskip=4pt,
  belowskip=4pt
}
\lstset{style=sqlstyle}

\title[SQL SELECT]{Introduction to SQL \texttt{SELECT} Statement}
\subtitle{Single-Table Querying with the \texttt{classicmodels} Sample Database}
\author{DBMS Course}
\institute{Smart Logistics and Supply Chain Management Lab}
\date{\today}

\newcommand{\kw}[1]{\textcolor{sqlblue}{\texttt{#1}}}
\newcommand{\minihead}[1]{\vspace{2pt}\textbf{#1}\vspace{2pt}}

\begin{document}

\begin{frame}
  \titlepage
\end{frame}

\begin{frame}{Learning Objectives}
After this lecture, students will be able to:
\begin{enumerate}
  \item Understand the purpose and role of the \kw{SELECT} statement in SQL.
  \item Retrieve all columns or specific columns from a table.
  \item Use column aliases and define computed columns in query results.
  \item Filter data using \kw{WHERE} and comparison operators.
  \item Sort results, limit returned rows, and eliminate duplicate values.
  \item Write practical single-table queries on the \texttt{classicmodels} database.
\end{enumerate}
\end{frame}

\begin{frame}[fragile]{What is SELECT Used For?}
\begin{block}{Core Concept}
\kw{SELECT} is used to retrieve data from one or more tables and present the result in a tabular format of rows and columns.
\end{block}
\begin{columns}[T]
\column{0.48\textwidth}
\minihead{Business Questions}
\begin{itemize}
  \item Which customers are located in the USA?
  \item Which products have an MSRP above 150?
  \item Which orders have been cancelled?
\end{itemize}
\column{0.48\textwidth}
\minihead{Answered in SQL}
\begin{lstlisting}
SELECT customerName, country
FROM customers
WHERE country = 'USA';
\end{lstlisting}
\end{columns}
\end{frame}

\begin{frame}{Sample Database: classicmodels}
\begin{itemize}
  \item \texttt{classicmodels} models the business operations of a retailer selling scale model cars, motorcycles, planes, and ships.
  \item In this lecture, we focus exclusively on single-table queries.
\end{itemize}
\vspace{4pt}
\begin{center}
\begin{tabular}{ll}
\toprule
\textbf{Table} & \textbf{Description} \\
\midrule
\texttt{products} & Product catalog and inventory data \\
\texttt{customers} & Customer profiles and contact details \\
\texttt{orders} & Order headers \\
\texttt{orderdetails} & Order line items \\
\texttt{payments} & Customer payment transactions \\
\texttt{employees} & Employee directory \\
\texttt{offices} & Regional sales office locations \\
\texttt{productlines} & Product lines and categories \\
\bottomrule
\end{tabular}
\end{center}
\end{frame}

\begin{frame}[fragile]{Setting Up the Practice Environment}
\begin{columns}[T]
\column{0.50\textwidth}
\minihead{Select the Database}
\begin{lstlisting}
USE classicmodels;
\end{lstlisting}
\minihead{List All Tables}
\begin{lstlisting}
SHOW TABLES;
\end{lstlisting}
\column{0.50\textwidth}
\minihead{Inspect Table Schema}
\begin{lstlisting}
DESC products;
DESC customers;
DESC orders;
\end{lstlisting}
\begin{alertblock}{Best Practice}
Before writing queries, always check the available columns and their data types.
\end{alertblock}
\end{columns}
\end{frame}

\begin{frame}[fragile]{Basic Syntax of SELECT}
\begin{block}{General Template}
\begin{lstlisting}
SELECT column1, column2, ...
FROM table_name;
\end{lstlisting}
\end{block}
\begin{center}
\begin{tabular}{ll}
\toprule
\textbf{Clause} & \textbf{Meaning} \\
\midrule
\kw{SELECT} & Specifies the columns to return in the result \\
\texttt{column1, column2, ...} & Comma-separated list of projected columns \\
\kw{FROM} & Specifies the table containing the data \\
\texttt{table\_name} & Target table name to query \\
\bottomrule
\end{tabular}
\end{center}
\end{frame}

\begin{frame}[fragile]{First Query Examples}
\begin{columns}[T]
\column{0.52\textwidth}
\minihead{Retrieve Customer Names}
\begin{lstlisting}
SELECT customerName
FROM customers;
\end{lstlisting}
\minihead{Retrieve Order Numbers}
\begin{lstlisting}
SELECT orderNumber
FROM orders;
\end{lstlisting}
\column{0.44\textwidth}
\begin{block}{Reading a Query}
\begin{itemize}
  \item \kw{SELECT}: which columns?
  \item \kw{FROM}: from which table?
  \item Semicolon (\texttt{;}) terminates the statement.
\end{itemize}
\end{block}
\end{columns}
\end{frame}

\begin{frame}[fragile]{Retrieving All Columns with SELECT *}
\begin{lstlisting}
SELECT *
FROM productlines;
\end{lstlisting}
\begin{columns}[T]
\column{0.48\textwidth}
\begin{block}{When is it Useful?}
\begin{itemize}
  \item Quickly inspecting table data.
  \item Exploring table schemas while learning SQL.
\end{itemize}
\end{block}
\column{0.48\textwidth}
\begin{alertblock}{Avoid in Production}
On large tables, \kw{SELECT *} can retrieve unnecessary columns, slowing down network transfer and wasting memory. In real applications, specify only required columns.
\end{alertblock}
\end{columns}
\end{frame}

\begin{frame}[fragile]{Selecting Specific Columns}
\begin{columns}[T]
\column{0.48\textwidth}
\minihead{Products}
\begin{lstlisting}
SELECT productCode, productName, productLine, MSRP
FROM products;
\end{lstlisting}
\minihead{Customers}
\begin{lstlisting}
SELECT customerNumber, customerName, city, country
FROM customers;
\end{lstlisting}
\column{0.48\textwidth}
\begin{block}{Rule of Thumb}
Project only columns necessary for the analysis or user interface.
\end{block}
\begin{exampleblock}{Question}
To build a customer directory grouped by country, which columns do we need?
\end{exampleblock}
\end{columns}
\end{frame}

\begin{frame}[fragile]{Aliasing Columns with AS}
\begin{lstlisting}
SELECT
    productCode AS product_code,
    productName AS product_name,
    productLine AS product_line,
    MSRP AS suggested_retail_price
FROM products;
\end{lstlisting}
\begin{itemize}
  \item Aliases provide clearer, user-friendly column headers in query output.
  \item The \kw{AS} keyword is optional, but using it explicitly improves readability.
\end{itemize}
\end{frame}

\begin{frame}[fragile]{Computed Columns in SELECT}
\begin{columns}[T]
\column{0.56\textwidth}
\begin{lstlisting}
SELECT
    productCode,
    productName,
    buyPrice,
    MSRP,
    MSRP - buyPrice AS profitEstimate
FROM products;
\end{lstlisting}
\column{0.40\textwidth}
\begin{block}{Key Points}
\begin{itemize}
  \item Calculations can be written directly in the \kw{SELECT} clause.
  \item Computed columns only exist in the query result set.
  \item The underlying table schema is not altered.
\end{itemize}
\end{block}
\end{columns}
\vspace{4pt}
\begin{lstlisting}
SELECT quantityOrdered * priceEach AS lineTotal
FROM orderdetails;
\end{lstlisting}
\end{frame}

\begin{frame}[fragile]{Filtering Data with WHERE}
\begin{block}{General Template}
\begin{lstlisting}
SELECT column1, column2
FROM table_name
WHERE condition;
\end{lstlisting}
\end{block}
\begin{columns}[T]
\column{0.48\textwidth}
\minihead{Motorcycles Category}
\begin{lstlisting}
SELECT productCode, productName, productLine
FROM products
WHERE productLine = 'Motorcycles';
\end{lstlisting}
\column{0.48\textwidth}
\minihead{Customers in USA}
\begin{lstlisting}
SELECT customerNumber, customerName, country
FROM customers
WHERE country = 'USA';
\end{lstlisting}
\end{columns}
\end{frame}

\begin{frame}{Common Comparison Operators}
\begin{center}
\begin{tabular}{lll}
\toprule
\textbf{Operator} & \textbf{Meaning} & \textbf{Example} \\
\midrule
\texttt{=} & Equal to & \texttt{country = 'USA'} \\
\texttt{<>} & Not equal to & \texttt{status <> 'Shipped'} \\
\texttt{>} & Greater than & \texttt{MSRP > 150} \\
\texttt{<} & Less than & \texttt{quantityInStock < 1000} \\
\texttt{>=} & Greater than or equal to & \texttt{creditLimit >= 100000} \\
\texttt{<=} & Less than or equal to & \texttt{buyPrice <= 50} \\
\bottomrule
\end{tabular}
\end{center}
\begin{exampleblock}{Remember}
The condition in \kw{WHERE} evaluates to TRUE or FALSE for each row.
\end{exampleblock}
\end{frame}

\begin{frame}[fragile]{Numeric Filtering Examples}
\begin{columns}[T]
\column{0.50\textwidth}
\minihead{MSRP Greater Than 150}
\begin{lstlisting}
SELECT productCode, productName, MSRP
FROM products
WHERE MSRP > 150;
\end{lstlisting}
\column{0.46\textwidth}
\minihead{Stock Quantity Below 1000}
\begin{lstlisting}
SELECT productCode, productName, quantityInStock
FROM products
WHERE quantityInStock < 1000;
\end{lstlisting}
\end{columns}
\vspace{5pt}
\begin{lstlisting}
SELECT customerNumber, customerName, creditLimit
FROM customers
WHERE creditLimit >= 100000;
\end{lstlisting}
\end{frame}

\begin{frame}[fragile]{Combining Conditions with AND}
\begin{block}{Logic}
\kw{AND} requires all conditions to evaluate to TRUE.
\end{block}
\begin{lstlisting}
SELECT productCode, productName, productLine, MSRP
FROM products
WHERE productLine = 'Classic Cars'
  AND MSRP > 100;
\end{lstlisting}
\begin{lstlisting}
SELECT customerNumber, customerName, country, creditLimit
FROM customers
WHERE country = 'USA'
  AND creditLimit > 100000;
\end{lstlisting}
\end{frame}

\begin{frame}[fragile]{Combining Conditions with OR}
\begin{block}{Logic}
\kw{OR} requires at least one condition to evaluate to TRUE.
\end{block}
\begin{columns}[T]
\column{0.50\textwidth}
\minihead{Customers in USA or France}
\begin{lstlisting}
SELECT customerNumber, customerName, country
FROM customers
WHERE country = 'USA'
   OR country = 'France';
\end{lstlisting}
\column{0.46\textwidth}
\minihead{Motorcycles or Planes}
\begin{lstlisting}
SELECT productCode, productName, productLine
FROM products
WHERE productLine = 'Motorcycles'
   OR productLine = 'Planes';
\end{lstlisting}
\end{columns}
\end{frame}

\begin{frame}[fragile]{Negating Conditions with NOT}
\begin{columns}[T]
\column{0.50\textwidth}
\minihead{Using NOT}
\begin{lstlisting}
SELECT customerNumber, customerName, country
FROM customers
WHERE NOT country = 'USA';
\end{lstlisting}
\column{0.46\textwidth}
\minihead{Equivalent Form}
\begin{lstlisting}
SELECT customerNumber, customerName, country
FROM customers
WHERE country <> 'USA';
\end{lstlisting}
\end{columns}
\vspace{5pt}
\begin{lstlisting}
SELECT orderNumber, orderDate, status
FROM orders
WHERE status <> 'Shipped';
\end{lstlisting}
\end{frame}

\begin{frame}[fragile]{Pattern Matching with LIKE}
\begin{block}{Concept}
\kw{LIKE} performs pattern matching on string values using wildcards.
\end{block}
\begin{center}
\begin{tabular}{ll}
\toprule
\textbf{Wildcard} & \textbf{Meaning} \\
\midrule
\texttt{\%} & Matches zero or more characters \\
\texttt{\_} & Matches exactly one character \\
\bottomrule
\end{tabular}
\end{center}
\begin{lstlisting}
SELECT customerName, city, country
FROM customers
WHERE customerName LIKE '%Mini%';
\end{lstlisting}
\begin{lstlisting}
SELECT productCode, productName
FROM products
WHERE productName LIKE '196%';
\end{lstlisting}
\end{frame}

\begin{frame}[fragile]{Filtering by a Set of Values with IN}
\begin{block}{Concept}
\kw{IN} provides a clean, concise shorthand for multiple \kw{OR} equality tests on the same column.
\end{block}
\begin{columns}[T]
\column{0.50\textwidth}
\minihead{Using Multiple OR}
\begin{lstlisting}
WHERE country = 'USA'
   OR country = 'France'
   OR country = 'Germany'
\end{lstlisting}
\column{0.46\textwidth}
\minihead{Using IN}
\begin{lstlisting}
WHERE country IN ('USA', 'France', 'Germany')
\end{lstlisting}
\end{columns}
\vspace{4pt}
\begin{lstlisting}
SELECT customerNumber, customerName, country
FROM customers
WHERE country IN ('USA', 'France', 'Germany');
\end{lstlisting}
\end{frame}

\begin{frame}[fragile]{Range Filtering with BETWEEN}
\begin{block}{Concept}
\kw{BETWEEN a AND b} tests whether a value falls within the range from \texttt{a} to \texttt{b}.
\end{block}
\begin{lstlisting}
SELECT productCode, productName, MSRP
FROM products
WHERE MSRP BETWEEN 100 AND 150;
\end{lstlisting}
\begin{lstlisting}
SELECT orderNumber, orderDate, status
FROM orders
WHERE orderDate BETWEEN '2005-01-01' AND '2005-12-31';
\end{lstlisting}
\begin{alertblock}{Note}
In SQL, \kw{BETWEEN} is inclusive of both boundary endpoints.
\end{alertblock}
\end{frame}

\begin{frame}[fragile]{Checking for NULL Values}
\begin{block}{What is NULL?}
\kw{NULL} represents unknown, missing, or inapplicable data. Never test for \kw{NULL} using \texttt{=}.
\end{block}
\begin{columns}[T]
\column{0.50\textwidth}
\minihead{Find Rows with NULL}
\begin{lstlisting}
SELECT customerNumber, customerName, salesRepEmployeeNumber
FROM customers
WHERE salesRepEmployeeNumber IS NULL;
\end{lstlisting}
\column{0.46\textwidth}
\minihead{Find Rows Without NULL}
\begin{lstlisting}
SELECT customerNumber, customerName, salesRepEmployeeNumber
FROM customers
WHERE salesRepEmployeeNumber IS NOT NULL;
\end{lstlisting}
\end{columns}
\end{frame}

\begin{frame}[fragile]{Eliminating Duplicates with DISTINCT}
\begin{block}{Concept}
\kw{DISTINCT} removes duplicates and returns only unique values in the query result.
\end{block}
\begin{columns}[T]
\column{0.50\textwidth}
\minihead{May Contain Duplicates}
\begin{lstlisting}
SELECT country
FROM customers;
\end{lstlisting}
\column{0.46\textwidth}
\minihead{Unique Values Only}
\begin{lstlisting}
SELECT DISTINCT country
FROM customers;
\end{lstlisting}
\end{columns}
\begin{exampleblock}{Business Question}
In which distinct countries does our company currently have customers?
\end{exampleblock}
\end{frame}

\begin{frame}[fragile]{Sorting Results with ORDER BY}
\begin{block}{General Template}
\begin{lstlisting}
SELECT column1, column2
FROM table_name
ORDER BY column1 ASC, column2 DESC;
\end{lstlisting}
\end{block}
\begin{columns}[T]
\column{0.50\textwidth}
\minihead{Ascending (ASC)}
\begin{lstlisting}
SELECT productCode, productName, MSRP
FROM products
ORDER BY MSRP ASC;
\end{lstlisting}
\column{0.46\textwidth}
\minihead{Descending (DESC)}
\begin{lstlisting}
SELECT productCode, productName, MSRP
FROM products
ORDER BY MSRP DESC;
\end{lstlisting}
\end{columns}
\end{frame}

\begin{frame}[fragile]{Restricting Row Count with LIMIT}
\begin{block}{Concept}
\kw{LIMIT} restricts the number of rows returned. In MySQL, it is widely used for top-$N$ queries and pagination.
\end{block}
\begin{columns}[T]
\column{0.50\textwidth}
\minihead{First 5 Products}
\begin{lstlisting}
SELECT productCode, productName, MSRP
FROM products
LIMIT 5;
\end{lstlisting}
\column{0.46\textwidth}
\minihead{Top 5 Most Expensive Products}
\begin{lstlisting}
SELECT productCode, productName, MSRP
FROM products
ORDER BY MSRP DESC
LIMIT 5;
\end{lstlisting}
\end{columns}
\end{frame}

\begin{frame}{Lexical Order vs. Logical Execution Order}
\begin{columns}[T]
\column{0.47\textwidth}
\begin{block}{Written (Lexical) Order}
\begin{enumerate}
  \item \kw{SELECT}
  \item \kw{FROM}
  \item \kw{WHERE}
  \item \kw{ORDER BY}
  \item \kw{LIMIT}
\end{enumerate}
\end{block}
\column{0.49\textwidth}
\begin{block}{Logical Execution Order}
\begin{enumerate}
  \item Locate the table (\kw{FROM})
  \item Filter rows (\kw{WHERE})
  \item Project columns (\kw{SELECT})
  \item Sort result rows (\kw{ORDER BY})
  \item Restrict count (\kw{LIMIT})
\end{enumerate}
\end{block}
\end{columns}
\end{frame}

\begin{frame}[fragile]{A Complete Example Query}
\begin{lstlisting}
SELECT
    productCode AS product_code,
    productName AS product_name,
    productLine AS product_line,
    MSRP AS suggested_price
FROM products
WHERE productLine IN ('Classic Cars', 'Motorcycles')
  AND MSRP BETWEEN 100 AND 200
ORDER BY MSRP DESC
LIMIT 10;
\end{lstlisting}
\begin{block}{Query Objective}
Retrieve the top 10 most expensive products in \texttt{Classic Cars} and \texttt{Motorcycles} whose suggested retail price is between 100 and 200.
\end{block}
\end{frame}

\begin{frame}{Common Pitfalls for Beginners}
\begin{center}
\begin{tabular}{>{\raggedright\arraybackslash}p{0.40\textwidth}>{\raggedright\arraybackslash}p{0.48\textwidth}}
\toprule
\textbf{Mistake} & \textbf{Fix} \\
\midrule
Missing commas between columns & Write \texttt{SELECT col1, col2, col3} \\
Missing \kw{FROM} clause & Always define the data source table \\
Unquoted string literals & Use single quotes: \texttt{country = 'USA'} \\
Comparing \kw{NULL} with \texttt{=} & Use \kw{IS NULL} or \kw{IS NOT NULL} \\
Overusing \kw{SELECT *} & Select only necessary columns \\
\bottomrule
\end{tabular}
\end{center}
\end{frame}

\begin{frame}[fragile]{Quick Exercise 1}
Write SQL queries for each requirement:
\begin{enumerate}
  \item Retrieve \texttt{customerNumber}, \texttt{customerName}, \texttt{city}, and \texttt{country} from \texttt{customers}.
  \item Retrieve customers located in \texttt{France} or \texttt{Germany}.
  \item Retrieve products with an \texttt{MSRP} greater than 150.
  \item Retrieve the 5 most recent orders based on \texttt{orderDate}.
\end{enumerate}
\pause
\begin{block}{Hint for Question 4}
\begin{lstlisting}
SELECT orderNumber, orderDate, status
FROM orders
ORDER BY orderDate DESC
LIMIT 5;
\end{lstlisting}
\end{block}
\end{frame}

\begin{frame}[fragile]{Quick Exercise 2}
Write SQL queries for each requirement:
\begin{enumerate}
  \item Retrieve the distinct list of countries from \texttt{customers}.
  \item Retrieve products whose name contains the substring \texttt{'Ford'}.
  \item Retrieve customers who do not have an assigned sales representative.
  \item Compute \texttt{lineTotal = quantityOrdered * priceEach} from \texttt{orderdetails}.
\end{enumerate}
\pause
\begin{block}{Hints / Solutions}
\begin{lstlisting}
SELECT DISTINCT country FROM customers;
SELECT * FROM products WHERE productName LIKE '%Ford%';
SELECT * FROM customers WHERE salesRepEmployeeNumber IS NULL;
\end{lstlisting}
\end{block}
\end{frame}

\begin{frame}{Summary}
\begin{itemize}
  \item \kw{SELECT} is the fundamental statement for querying data in relational databases.
  \item \kw{FROM} specifies the source table for the query.
  \item \kw{WHERE} filters rows based on logical conditions.
  \item \kw{AS} assigns aliases to columns; expressions in \kw{SELECT} produce computed columns.
  \item \kw{DISTINCT}, \kw{ORDER BY}, and \kw{LIMIT} clean, sort, and paginate query results.
  \item When drafting queries, think in business questions: \textit{which columns, from which table, under which conditions?}
\end{itemize}
\end{frame}

\begin{frame}{References}
\small
\begin{itemize}
  \item DBMS Tutorial: \textit{Basic SELECT Statement with classicmodels}.\\
  {\tiny\url{https://vudmvn.github.io/dbms-course/lectures/MySQL/querying-data/select/select-statement-1/}}
  \item GeeksforGeeks: \textit{SQL SELECT Query}.\\
  {\tiny\url{https://www.geeksforgeeks.org/sql/sql-select-query/}}
\end{itemize}
\end{frame}

\end{document}
